%%
%% fall-lecture16.tex
%% 
%% Made by Alex Nelson <pqnelson@gmail.com>
%% Login   <alex@lisp>
%% 
%% Started on  2026-11-08T11:29:38-0800
%% Last update 2026-11-08T11:29:38-0800
%% 

\lecture{}

\begin{recall}
We gave axioms for Steenrod square operation 
\begin{equation}
\Sq^{i}\colon H^{k}(X;\ZZ_{2})\to H^{k+i}(X;\ZZ_{2})
\end{equation}
is $\ZZ_{2}$-linear and satisfies:
\begin{enumerate}
\item naturality: $\Sq^{i}(f^{*}\alpha)=f^{*}\Sq^{i}(\alpha)$ [i.e.,
  commutes with pullbacks]
\item Cartan's formula: $\Sq^{i}(\alpha\smile\beta)=\sum_{j}\Sq^{j}(\alpha)\smile\Sq^{i-j}(\beta)$
\item When $H^{n}(X;\ZZ_{2})$ is the top cohomology group, for any
  $\alpha\in H^{n}(X;\ZZ_{2})$ we have $\Sq^{n}(\alpha)=\alpha\smile\alpha$.
  For any $i>n$, for any $\alpha\in H^{n}(X;\ZZ_{2})$, we have $\Sq^{i}(\alpha)=0$.
\item $\Sq^{0}(\alpha)=\alpha$.
\end{enumerate}
These are the relevant properties we will need.
\end{recall}

\begin{node}
We define the Stiefel--Whitney classes using the Thom isomorphism and
Steenrod square operation. For any bundle $\xi\colon E\xrightarrow{\pi}B$
we have the Thom class $U\in H^{n}(E,E_{0};\ZZ_{2})$ and the Thom
isomorphism
\begin{align}
  \varphi\colon & H^{i}(B;\ZZ_{2})\to H^{i+n}(E,E_{0};\ZZ_{2})\\
  &\beta\mapsto(\pi^{*}\beta)\smile U\nonumber
\end{align}
Then we define the Stiefel--Whitney class on $\xi$ as:
\begin{equation}
\boxed{w_{i}(\xi)=\varphi^{-1}(\Sq^{i}(U))=\varphi^{-1}(\Sq^{i}(\varphi(1))).}
\end{equation}
\end{node}

\begin{remark}[History]
An abbreviated history behind this proposed definition for the
Stiefel--Whitney class may be helpful.
\begin{enumerate}
\item In 1935, Stiefel and Whitney independently introduced the
  Stiefel--Whitney class from the obstruction theory
  point-of-view. Here $w_{n-k+1}(\xi)$ is the ``obstruction mod~2'' of
  the existence of a $k$-nowhere dependent sections over the
  $(n-k+1)$-skeleton of $B$ [where ``obstruction mod~2'' means
    $w_{n-k+1}\in H^{n-k+1}(B;\ZZ_{2})$].
\item In 1947, Steenrod introduced $\Sq^{i}$
\item In 1951, Thom (in his PhD thesis) discovered the formula for
  embedded submanifolds (for normal bundles). His advisor, Henri
  Cartan, observed it could be extended to general vector bundles.
\end{enumerate}
How did Thom do it/ Let's look at an $m$-manifold $M$ and embed it in
a Riemannian $(m+r)$-manifold $N$. Let us look at $E$ the rank-$r$
normal bundle of $M$. Then $\varphi(w_{r}(E))=\Sq^{r}(U)$ where $U\in H^{r}(E,E_{0};\ZZ_{2})$
is the Thom class of $E$. The $w_{r}(E)$ is from the original
definition via obstruction theory.

See Milnor and Stasheff~\cite[Ch~12]{milnor1974characteristic} for
more details about the history.
\end{remark}

\begin{observation}[Thom?]
We can realize
\begin{subequations}
  \begin{align}
w_{r}(E) &= e_{2}(E)\quad\mbox{(= Euler class mod 2)}\\
&= PD_{M}[Z]
  \end{align}
\end{subequations}
where $PD_{M}$ is the Poincar\'{e} duality relating cohomology with
homology, and $Z\subset M$ is a submanifold (so $[Z]$ is its
fundamental class). We have a natural zero section
\begin{equation}
z\colon M\to E
\end{equation}
and we can perturb it ``a little bit'' to obtain a general section
\begin{equation}
s\colon M\to E.
\end{equation}
We need $s$ to span the tangent space at the intersection with $z$
[spanning the transversal direction]. We can construct the
$(m-r)$-manifold
\begin{equation}
Z=\{x\in M\mid s(x)=0\}\subset M
\end{equation}
which is a submanifold of $M$. Then
\begin{equation}
s\colon(M,M\setminus Z)\to(E,E\setminus z(M))
\end{equation}
is how we relate $w_{r}$ to Poincar\'{e} duality $PD_{M}[Z]$.

The inclusion
\begin{equation}
j\colon H^{r}(E,E\setminus z(M);\ZZ_{2})\to H^{r}(E;\ZZ_{2})
\end{equation}
maps $U\mapsto j(U)$. We see the Euler class mod~2 is the pullback
$e_{2}(E)=s^{*}j(U)\in H^{r}(M;\ZZ_{2})$. The Poincar\'{e}-duality is
$PD_{M}[Z]=s^{*}j(U)\in H^{r}(M;\ZZ_{2})$, and $w_{r}(E)=PD_{M}[Z]$ is
the key step.

We see
\begin{subequations}
  \begin{align}
\varphi(w_{r}(E)) &= \pi^{*}(w_{r}(E))\smile U\\
&= \pi^{*}z^{*}j(U)\smile U\\
&=(z\circ\pi)^{*}j(U)\smile U\quad\mbox{but $z\circ\pi\homotopic\id$ homotopic}\\
&=j(U)\smile U\\
&= U\smile U.
  \end{align}
\end{subequations}
Thom asked, ``What about the $(r-1)$ case?''
\end{observation}

\subsection{Existence}

\begin{claim}
$w_{i}(\xi)=\varphi^{-1}(\Sq^{i}(U))=\varphi^{-1}(\Sq^{i}(\varphi(1)))$
  satisfies all 4 axioms for the Stiefel--Whitney class.
\end{claim}

\begin{Proof}[Proof (Axiom 1)]
We see $w_{i}(\xi)\in H^{i}(B;\ZZ_{2})$ since $\Sq^{i}(U)\in H^{n+1}(E,E_{0};\ZZ_{2})$
and $\varphi^{-1}\Sq^{i}(U)\in H^{i}(B;\ZZ_{2}$.

We see
\begin{subequations}
  \begin{align}
w_{0}(\xi) &= \varphi^{-1}(\Sq^{0}(\varphi(1)))\\
&= \varphi^{-1}\circ\varphi(1)\\
&= 1.
  \end{align}
\end{subequations}
For $i>n$, we see
\begin{subequations}
  \begin{align}
w_{i}(\xi) &= \varphi^{-1}(\Sq^{i}(U))\\
&= \varphi^{-1}(0)\\
&= 0.
  \end{align}
\end{subequations}
Hence this satisfies Axiom~1.
\end{Proof}

\begin{Proof}[Proof (Axiom 2)]
Let $\xi\colon E\xrightarrow{\pi}B$, $f\colon A\to B$ be continuous.
We will show $w_{i}(f^{*}\xi)=f^{*}w_{i}(\xi)$. We extend $f$ by the
following commutative diagram
\begin{equation}
\vcenter{\xymatrix{\widetilde{E}\ar[r]^{\widetilde{f}}\ar[d]_{\pi}&E\ar[d]^{\pi}\\
A\ar[r]^{f} & B}}
\end{equation}
where $\widetilde{f}(a,x)=x$. So this means $\pi\circ\widetilde{f}=f\circ\pi$.
Then if we want to prove the claim, what do we need to do?

Well, $f^{*}w_{i}(\xi)=f^{*}\varphi^{-1}_{\xi}(\Sq^{i}(U))$ and
$w_{i}(f^{*}E)=\varphi^{-1}_{f^{*}\xi}\Sq^{i}(\widetilde{U})$ where
$\widetilde{U}\in H^{n}(\widetilde{E},\widetilde{E}_{0};\ZZ_{2})$.
What is the relation between $U$ and $\widetilde{U}$? Well, abusing
notation and introducing a new map confusingly called
\begin{equation}
\widetilde{f}\colon(\widetilde{E},\widetilde{E}_{0})\to(E,E_{0})
\end{equation}
then $(\widetilde{f})^{*}U=\widetilde{U}$. Then we plug this into the
left-hand side of the claim:
\begin{subequations}
  \begin{align}
\varphi^{-1}_{f^{*}\xi}(\Sq^{i}\widetilde{U})
&=\varphi^{-1}_{f^{*}\xi}(\Sq^{i}((\widetilde{f})^{*}U))\\
&=\varphi^{-1}_{f^{*}\xi}((\widetilde{f})^{*}\Sq^{i}(U)).
  \end{align}
\end{subequations}
Then by the naturality of the Thom isomorphism,
\begin{equation}
f^{*}\varphi^{-1}_{\xi}=\varphi^{-1}_{f^{*}\xi}(\widetilde{f})^{*}\iff\varphi^{-1}_{f^{*}\xi}f^{*}=(\widetilde{f})^{*}\varphi_{\xi}.
\end{equation}
Let $\alpha\in H^{i}(B;\ZZ_{2})$, then the right-hand side of the
claim becomes
\begin{subequations}
  \begin{align}
(\widetilde{f})^{*}\varphi_{\xi}(\alpha)
&=(\widetilde{f})^{*}(\pi^{*}\alpha\smile U)\\
&=(\pi\circ\widetilde{f})^{*}(\alpha)\smile(\widetilde{f})^{*}U\\
&=(f\circ\pi)^{*}(\alpha)\smile(\widetilde{f})^{*}U,
  \end{align}
\end{subequations}
and the left-hand side transforms further into
\begin{subequations}
  \begin{align}
LHS &= \pi^{*}(f^{*}\alpha)\smile(\widetilde{f})^{*}U\\
&= (f\circ\pi)^{*}\alpha\smile(\widetilde{f})^{*}U,
  \end{align}
\end{subequations}
from which it obviously follows the claim holds.
\end{Proof}

\begin{Proof}[Proof (Axiom 3)]
The Whitney sum property $w(\xi\oplus\eta)=w(\xi)w(\eta)$.
We can define the following notation
\begin{equation}
\Sq(-)=\Sq^{0}(-)+\Sq^{1}(-)+\dots+\Sq^{n}(-)
\end{equation}
Then
\begin{equation}
w(\xi)=\varphi^{-1}\Sq(U),
\end{equation}
where $U$ is the Thom class.

Recall, $\Delta^{*}(\xi\times\eta)=\xi\oplus\eta$ when we pullback the
product bundle along the diagonal map, we form the Whitney sum of bundles.

It suffices to show $w(\xi\times\eta)=w(\xi)w(\eta)$.
We see that
\begin{equation}
LHS = \varphi_{\xi\times\eta}^{-1}\Sq(U_{\xi\times\eta}).
\end{equation}
What is the relation between the various Thom classes
$U_{\xi\times\eta}$, $U_{\xi}$, $U_{\eta}$? Let $U\in H^{m}(E,E_{0};\ZZ_{2})$
be the Thom class of the rank-$m$ bundle $\xi\colon E\to B$, let
$U'\in H^{n}(E',E'_{0};\ZZ_{2})$ be the Thom class of the rank-$n$
bundle $\eta\colon E'\to B$, then
\begin{equation}
U_{\xi\times\eta}=U\times U'\in H^{m+n}(E\times E',\underbrace{(E\times
E_{0}')\cup(E_{0}\times E')}_{=(E\times E')_{0}};\ZZ_{2}).
\end{equation}
We see $U\times U'$ is well-defined. We need to prove $U\times U'$ is
a Thom class.

Let $(x,y)\in B\times B$, then look at the inclusion mapping
\begin{equation}
i\colon(F_{(x,y)},F_{(x,y),0})\to(E\times E',(E\times E')_{0}),
\end{equation}
so
\begin{equation}
i^{*}(U\times U') = i^{*}U_{x}\times i^{*}U_{y}
\end{equation}
where $i^{*}U_{x}\in H^{m}(F_{x},(F_{x})_{0};\ZZ_{2})$ and
$i^{*}U_{y}\in H^{n}(F_{y},(F_{y})_{0};\ZZ_{2})$.

We see that the Thom isomorphism
\begin{equation}
\varphi_{\xi\times\eta}(\alpha\times\beta)=\varphi_{\xi}(\alpha)\times\varphi_{\eta}(\beta),
\end{equation}
by similar reasoning.

OK, so
\begin{subequations}
  \begin{align}
w(\xi\times\eta)
&=\varphi_{\xi\times\eta}^{-1}(\Sq(U_{\xi\times\eta}))\\
&=\varphi_{\xi\times\eta}^{-1}(\Sq(U\times U'))\\
&=\varphi_{\xi\times\eta}^{-1}(\Sq(U)\Sq(U'))\quad\mbox{by Cartan's formula}\\
&=\varphi_{\xi}^{-1}(\Sq(U))\varphi_{\eta}^{-1}(\Sq(U'))\\
&=w(\xi)w(\eta).
  \end{align}
\end{subequations}
Hence this construction satisfies Axiom~3.
\end{Proof}

Next week we will hopefully start Complex bundles. Next time, we will
finish proving existence by showing our construction satisfies Axiom~4.